\newpage
\section{Acta Número 21}
\label{sec:acta21}

\noindent \textbf{Fecha de la sesión:}\par
Jueves, 14 de Mayo de 2026

\vspace*{0.5cm}
\noindent \textbf{Datos de la Sesión:}
\vspace{-0.2cm}
\begin{itemize}[noitemsep]
\item \textbf{Modalidad:} Virtual - Google Meet.
\item \textbf{Hora de Inicio:} 18:00
\item \textbf{Hora de Fin:} 20:00
\end{itemize}

\vspace*{0.5cm}
\noindent \textbf{Orden del Día:}
\vspace{-0.2cm}
\begin{itemize}[noitemsep]
\item Control de Asistencia.
\item Revisión de consultas técnicas sobre el desarrollo técnico de las capas del sistema y socialización de modificaciones en la base de datos.
\item Planificación, análisis y viabilidad del despliegue en entornos oficiales universitarios (UMSS - Webtis).
\item Evaluación y determinación de alternativas de infraestructura externa para el despliegue del sistema.
\item Firmas y Aprobación de los Presentes.
\nobreak\end{itemize}

\vspace*{0.5cm}
\noindent \textbf{Socios Fundadores Presentes:}
\vspace{-0.2cm}
\begin{enumerate}[noitemsep]
\item Pardo Pozo Juan Diego (Jefe de Proyecto).
\item Arnez Jaimes Mauricio.
\item Quispe Flores Camila Belen.
\item Ariñez Bautista Jim Gabriel.
\item Medina Rodríguez Mateo Jhoustin.
\item Molina Maldonado Tomas Manuel.
\item Gutierrez Guzmán Angheline.
\end{enumerate}

\newpage
\subsection{Desarrollo del Acta}

\noindent \textbf{Control de Asistencia del Team}\par
Se procedió a la toma de asistencia a la hora acordada, corroborando la presencia de los 7 socios fundadores de la grupo-empresa Byte Busters S.R.L. Al contar con el quórum y la atención de todos los involucrados, se dio inicio formal a la sesión.

\vspace*{0.5cm}
\noindent \textbf{Revisión de Consultas Técnicas e Integraciones de Base de Datos}\par
El equipo se reunió con el propósito de centralizar y absolver dudas vinculadas al desarrollo activo de la plataforma. Durante este bloque, se abrieron espacios de discusión enfocados en las capas de \textit{Backend}, \textit{Frontend} y el modelado de la \textit{Base de Datos}. Tras un intercambio colaborativo, se constató con agrado que no existían dudas pendientes en ninguna de las áreas, demostrando la sólida comprensión del equipo sobre la arquitectura implementada. Asimismo, se realizó una explicación técnica pormenorizada sobre las nuevas adiciones e integraciones estructurales en la base de datos, garantizando que todos los desarrolladores estén plenamente alineados para las fases consecutivas.

\vspace*{0.5cm}
\noindent \textbf{Análisis de Viabilidad de Despliegue en Servidores Oficiales}\par
Como segundo punto crucial de la agenda, se debatió a profundidad la estrategia de despliegue en los servidores institucionales de la Universidad Mayor de San Simón (UMSS - Webtis). La sesión se orientó a coordinar y prever soluciones operativas ante los retos de la infraestructura local. Con el fin de resguardar el correcto y fluido funcionamiento del sistema y, fundamentalmente, asegurar el cumplimiento impecable del marco contractual establecido con la Consultora TIS, el equipo evaluó minuciosamente los escenarios de riesgo y los requerimientos de disponibilidad técnica del software.

\vspace*{0.5cm}
\noindent \textbf{Determinación de Infraestructura Externa Alternativa}\par
Como resolución definitiva a las necesidades de despliegue abordadas, se llegó al consenso unánime de proyectar e implementar la puesta en producción del sistema utilizando servidores de entornos externos a la universidad. Esta decisión estratégica busca mitigar cualquier limitación técnica imprevista y proveer un entorno de alta disponibilidad. En este sentido, se acordó adoptar e integrar formalmente las plataformas tecnológicas de \textbf{Supabase} y \textbf{Vercel} como las principales alternativas viables para la gestión de datos y el hospedaje de la aplicación, asegurando un estándar óptimo de rendimiento y estabilidad para los usuarios finales.

\vspace*{0.5cm}
\noindent \textbf{Aprobación Oficial}\par
Al consolidar una claridad técnica absoluta en todos los frentes de trabajo, trazar una ruta de despliegue externa altamente eficiente que blinda los acuerdos con la Consultora TIS, y contar con el respaldo unánime de los integrantes, el equipo manifiesta su total conformidad. Se da por aprobada la sesión y se procede con la firma de los presentes.

\newpage
\noindent \textbf{Firmas y Aprobación de los Presentes}
\begin{center}
    \noindent
    \setlength{\tabcolsep}{0pt}
    \begin{tabular}{ccc}
        \espacioFirma{assets/img/firmas/mauricio.jpg}{Mauricio Jaimes Arnez}{13751221} & 
        \espacioFirma{assets/img/firmas/camila.jpg}{Camila Belen Quispe Flores}{13097244} &
        \espacioFirma{assets/img/firmas/jim.jpg}{Jim Gabriel Ariñez Bautista}{9517642} \\[1.0cm]
        \espacioFirma{assets/img/firmas/mateo.jpg}{Mateo Medina Rodríguez}{9453809} &
        \espacioFirma{assets/img/firmas/tomas.jpg}{Tomas Molina Maldonado}{8051701} &
        \espacioFirma{assets/img/firmas/angheline.jpg}{Angheline Gutierrez Guzmán}{13751815} \\[1.0cm]
    \end{tabular}
    \vspace{0.2cm} 
    \begin{tabular}{cc}
        \espacioFirma{assets/img/firmas/diego.jpg}{Juan Diego Pardo Pozo}{12747783}
    \end{tabular}
\end{center}

\vspace{0.5cm}
\noindent \textbf{Fecha de última revisión:} Jueves, 14 de Mayo de 2026